\chapter{Formalismo de Campo 6D y Tiempo Vectorial}
Evolucionamos de la red discreta a un campo escalar informacional $\Phi$ definido sobre un espacio de seis dimensiones: tres coordenadas espaciales $(x,y,z)$ y tres coordenadas temporales $(t_x, t_y, t_z)$. El tiempo ya no es una línea unidimensional, sino una estructura vectorial que captura diferentes aspectos del devenir.

\section{Las tres flechas del tiempo}
\begin{itemize}
    \item $t_1$: \textbf{Tiempo causal} – la flecha termodinámica asociada al aumento de entropía.
    \item $t_2$: \textbf{Tiempo informacional} – el grado de integración de la información (medida por la complejidad efectiva).
    \item $t_3$: \textbf{Tiempo estructural} – el cambio en la geometría del espacio y las relaciones topológicas.
\end{itemize}

\section{El Laplaciano 6D Optimizado}
La difusión de información en el tejido de coherencia se modela mediante un operador de Laplace extendido, implementable en arquitecturas GPU:
\begin{equation}
\nabla^2_6 \Phi = \sum_{i=1}^{3} \frac{\partial^2 \Phi}{\partial x_i^2} + \sum_{\alpha=1}^{3} \frac{\partial^2 \Phi}{\partial c_\alpha^2}
\end{equation}
donde $c_\alpha$ representan las dimensiones temporales.

\section{Dinámica de Cristalización}
La ecuación de evolución integra memoria exponencial e interacciones no lineales:
\begin{equation}
\frac{\partial^2 \Phi}{\partial t^2} = c^2 \nabla^2_6 \Phi - \alpha \Phi - \beta \Phi^3 - \gamma \frac{\partial \Phi}{\partial t} + \eta \mathcal{M}[\Phi]
\end{equation}
Aquí $\gamma$ es la disipación que elimina cristales de tiempo parásitos, y $\mathcal{M}[\Phi]$ es un término de memoria informacional que depende de la historia reciente del campo. Las soluciones estacionarias corresponden a configuraciones de máxima coherencia (cristales espacio-temporales).
